% ============================================================================
% Fichier  : partie1/P01_C04_panorama_referentiels.tex
% Titre    : Panorama et Articulation des Référentiels
% Auteur   : MINKA MI NGUIDJOÏ Thierry Emmanuel
% Version  : 1.0 -- octobre 2026
% ============================================================================

\chapter{Panorama et Articulation des Référentiels}
\label{chap:referentiels}
\chaptermark{Panorama des référentiels}

\epigraph{%
    « Un référentiel ne pense pas à votre place.\\
    Il vous évite seulement d'oublier ce que d'autres ont appris à leurs
    dépens. »%
}{Adage des praticiens de la gouvernance}

\niveauChapitre{Fondamental}{%
    Maîtrise du Chapitre~\ref{chap:alignement} (alignement, droits de
    décision, trois finalités de la gouvernance).%
}

\begin{boxobjectifs}
\begin{itemize}
    \item Situer COBIT, ITIL, ISO/IEC~38500, ISO/IEC~20000, ISO/IEC~27001,
          COSO, CMMI, TOGAF et PMBOK.
    \item Distinguer cadre, norme certifiable et bonnes pratiques.
    \item Associer chaque référentiel au besoin qu'il couvre.
    \item Construire une carte d'articulation des référentiels pour une
          organisation donnée.
\end{itemize}
\end{boxobjectifs}

\bigskip

% ============================================================================
\section*{Question directrice : faut-il tout appliquer~?}
% ============================================================================

Après la panne, le conseil d'administration de la BANQUE SANAGA reçoit en
trois semaines quatre propositions de cabinets de conseil. La première
propose « une mise en conformité COBIT », la deuxième « un déploiement
ITIL », la troisième « une certification ISO~27001 », la quatrième « un
programme de maturité CMMI ». Chacune promet de régler le problème. Le
président du conseil demande à la nouvelle responsable GRC~: « Faut-il
tout faire~? Dans quel ordre~? Et ces référentiels se contredisent-ils~? »

Ce chapitre répond à ces questions. Il commence par distinguer les natures
de référentiels (section~\ref{sec:c04-natures}), présente ensuite les
principaux d'entre eux, regroupés par besoin
(sections~\ref{sec:c04-gouvernance} à~\ref{sec:c04-projets}), puis montre
comment les articuler en une carte cohérente
(section~\ref{sec:c04-carte}). La réponse courte est~: \textbf{non, il ne
faut pas tout appliquer}, et les référentiels sérieux ne se contredisent
pas~: ils se répartissent le travail.

\separateur

% ============================================================================
\section{Trois Natures de Référentiels}
\label{sec:c04-natures}
% ============================================================================

\subsection{Cadre, norme, bonnes pratiques}

Le mot « référentiel » recouvre des documents de nature très différente.
Les confondre conduit à des promesses impossibles, comme celle d'une
« certification COBIT » d'entreprise, qui n'existe pas.

\begin{boxdef}
\termecle{Cadre de référence}\index{cadre de référence}
(\textit{framework})~: ensemble structuré de principes, de concepts et
d'objectifs qui organise un domaine et guide sa mise en œuvre, sans
imposer d'exigences vérifiables par un tiers. Exemples~: COBIT, COSO.

\medskip
\termecle{Norme certifiable}\index{norme certifiable}
(\textit{certifiable standard})~: norme qui formule des \emph{exigences}
(le verbe « doit »), dont le respect peut être audité par un organisme de
certification accrédité, qui délivre alors un certificat à
l'organisation. Exemples~: ISO/IEC~27001, ISO/IEC~20000-1, ISO~22301.

\medskip
\termecle{Recueil de bonnes pratiques}\index{bonnes pratiques}
(\textit{body of knowledge, guidance})~: ensemble de pratiques reconnues,
présentées comme des recommandations (« il convient de »), que
l'organisation adapte librement. Exemples~: ITIL, PMBOK, ISO/IEC~27002,
ISO~31000.
\end{boxdef}

\subsection{Certifier une organisation, certifier une personne}

La seconde confusion porte sur ce qui est certifié. Une organisation peut
être certifiée conforme à une norme d'exigences. Une personne peut être
certifiée compétente dans un référentiel. Les deux n'ont pas la même
portée.

\begin{tableaucours}{L{3.0cm}L{4.4cm}L{4.8cm}}{%
    \tableauHeader{Référentiel} &
    \tableauHeader{Certification de l'organisation} &
    \tableauHeader{Certification des personnes (exemples)}
}
COBIT & Non (évaluation de capacité possible) & COBIT Foundation, CGEIT,
CISA (ISACA) \\
ITIL & Non (voir ISO/IEC~20000-1) & ITIL Foundation et niveaux supérieurs
(PeopleCert) \\
ISO/IEC 20000-1 & Oui & Auditeur ou responsable de mise en œuvre \\
ISO/IEC 27001 & Oui & Auditeur ou responsable de mise en œuvre \\
ISO/IEC 38500 & Non & Formations spécialisées \\
COSO & Non & Certificats de l'IIA et d'autres organismes \\
CMMI & Évaluation de niveau de maturité (pas une certification ISO) &
Évaluateurs agréés \\
TOGAF & Non & TOGAF (The Open Group) \\
PMBOK & Non & PMP et autres (PMI) \\
\end{tableaucours}
\vspace{-0.8em}
\begin{center}{\small\itshape Tableau~4.1. Ce qui se certifie, et pour qui}\end{center}

\begin{boxalert}
\textbf{Violation typique.} Une organisation obtient la certification
ISO~27001 et en conclut que sa gouvernance des SI est « conforme ». Or le
périmètre du certificat peut ne couvrir qu'un centre de données, et la
norme ne traite que du management de la sécurité de l'information~: elle
ne dit rien de la valeur des investissements ni de la qualité des
services.

\smallskip
\textbf{Règle~:} lire toujours le \emph{périmètre} d'un certificat, et ne
jamais confondre la certification d'un système de management avec une
gouvernance des SI d'ensemble.
\end{boxalert}

\subsection{Le cadre légal camerounais}

Une certification volontaire ne remplace pas les obligations légales.

\begin{boxjuris}
\textbf{\loicam{n°~2010/012} du 21~décembre~2010 relative à la
cybersécurité et la cybercriminalité au Cameroun, article~11} (version
anglaise officielle)~:

\medskip
\textit{« The following activities may be subject to authorization:
\textelp{} the provision of security auditing, security programs editing,
and other authorized security services to the public. »}

\medskip
Soit~: la fourniture au public de services d'audit de sécurité peut être
soumise à autorisation.

\medskip
\textbf{Portée.} Au Cameroun, l'audit de sécurité des systèmes
d'information n'est pas une activité libre au même titre que le conseil
en organisation. Les articles~13 et~32 de la même loi
(Chapitre~\ref{chap:prologue}) rendent par ailleurs cet audit obligatoire
pour certains systèmes. Un certificat ISO/IEC~27001, délivré par un
organisme étranger, ne dispense donc pas des audits prévus par la loi~:
les deux démarches coexistent, et une bonne gouvernance les coordonne. Le
Chapitre~\ref{chap:droit} précisera les textes d'application.
\end{boxjuris}

% ============================================================================
\section{Les Référentiels de la Gouvernance}
\label{sec:c04-gouvernance}
% ============================================================================

\subsection{COBIT : le cadre de gouvernance et de management}

COBIT est publié par l'ISACA, association internationale des
professionnels de l'audit et de la gouvernance des SI. Né en 1996 comme
un référentiel de \emph{contrôle} à l'usage des auditeurs, il est devenu
un cadre de \emph{gouvernance et de management} de l'ensemble de
l'information et des technologies de l'organisation. Sa version
actuelle, COBIT~2019~\cite{isaca2018cobitintro}, compte 40~objectifs
répartis entre un domaine de gouvernance (EDM) et quatre domaines de
management (APO, BAI, DSS, MEA). Elle fait l'objet de la
Partie~\ref{part:cobit}.

\textbf{Sa force}~: c'est le seul référentiel qui couvre de bout en bout
le domaine, de la décision du conseil au contrôle de l'exploitation. Il
sert de « référentiel parapluie » sous lequel les autres viennent se
ranger. \textbf{Sa limite}~: il dit \emph{quoi} maîtriser, rarement
\emph{comment} le faire en détail.

\subsection{ISO/IEC 38500 : les principes pour l'organe de gouvernance}

La norme ISO/IEC~38500 s'adresse directement à l'organe de gouvernance
(conseil d'administration, direction). Sa troisième édition date de
février~2024~\cite{iso38500_2024}. Courte, elle ne décrit pas de
processus, mais des principes et un modèle de gouvernance fondé sur trois
tâches~: \emph{évaluer}, \emph{diriger}, \emph{surveiller}. Elle n'est
pas certifiable. Elle est étudiée au Chapitre~\ref{chap:normes}.

\subsection{COSO : le contrôle interne et la gestion des risques}

Le Committee of Sponsoring Organizations of the Treadway Commission
(\sigle{COSO}) publie deux cadres de référence. Le premier, sur le
\emph{contrôle interne} (1992, mis à jour en 2013), est la référence des
auditeurs et des régulateurs financiers, notamment pour la loi
Sarbanes-Oxley. Le second, sur la \emph{gestion des risques de
l'entreprise} (2004, refondu en 2017 sous le titre \emph{Enterprise Risk
Management: Integrating with Strategy and Performance}), relie le risque à
la stratégie et à la performance~\cite{coso2017erm}. COSO n'est pas
propre à l'informatique, mais la gouvernance des SI s'y inscrit~: les
contrôles généraux informatiques font partie du contrôle interne.

% ============================================================================
\section{Les Référentiels des Services et de la Sécurité}
\label{sec:c04-services}
% ============================================================================

\subsection{ITIL : les bonnes pratiques de la gestion des services}

ITIL est né à la fin des années 1980 au sein de l'administration
britannique, qui cherchait à rationaliser la gestion de ses services
informatiques. Ses versions successives l'ont fait passer d'un catalogue
de processus (version~2) à un cycle de vie des services (version~3, 2007),
puis à un système de valeur des services organisé en pratiques
(ITIL~4, 2019). Il est aujourd'hui la propriété de PeopleCert, qui a
lancé en 2026 \emph{ITIL version~5}~\cite{peoplecert2026itil5}.

\textbf{Sa force}~: une description détaillée et éprouvée de la manière de
gérer les services au quotidien (incidents, changements, niveaux de
service\ldots). \textbf{Sa limite}~: il ne traite la gouvernance qu'à la
marge, et ce n'est pas une norme~: on ne « certifie » pas une organisation
ITIL.

\begin{boiteRemarque}{De ITIL~4 à ITIL version~5}
ITIL version~5 a été annoncée par PeopleCert en janvier~2026, et ses
modules de formation sont publiés progressivement depuis. Selon la
présentation de l'éditeur, elle passe d'une approche centrée sur la
gestion des services à une approche intégrée des \emph{produits et des
services numériques}, couvre leur cycle de vie de bout en bout, et se
veut conçue d'emblée pour un environnement où l'intelligence artificielle
est omniprésente~\cite{peoplecert2026itil5}. Elle conserve les sept
principes directeurs et l'essentiel des pratiques d'ITIL~4. ITIL~4 reste
commercialisé et ses certifications restent reconnues pendant la
transition.

\smallskip
Ce cours s'appuie sur \textbf{ITIL~4} dans la Partie~\ref{part:pratique}~:
c'est la version documentée par des sources publiques, et celle que
pratiquent aujourd'hui les organisations. Le Chapitre~\ref{chap:itil-svs}
situera ITIL version~5 par rapport à elle.
\end{boiteRemarque}

\subsection{ISO/IEC 20000-1 : la norme du management des services}

ISO/IEC~20000-1 formule les exigences d'un \emph{système de management des
services}. Elle est certifiable, et c'est par elle qu'une organisation
peut faire reconnaître par un tiers la qualité de sa gestion des services.
Ses exigences sont très proches des pratiques d'ITIL, dont elle est
historiquement issue~: ITIL dit \emph{comment faire}, ISO/IEC~20000-1 dit
\emph{ce qui doit être démontré}.

\subsection{ISO/IEC 27001 et 27002 : la sécurité de l'information}

ISO/IEC~27001 formule les exigences d'un \emph{système de management de la
sécurité de l'information} (\sigle{SMSI}) et est certifiable. ISO/IEC~27002
détaille les mesures de sécurité qui peuvent y être mises en œuvre. Leur
édition en vigueur date de 2022. Elles s'appuient sur une appréciation des
risques, que précise ISO/IEC~27005. La famille ISO/IEC~27000 sera reprise
au Chapitre~\ref{chap:normes}.

\subsection{ISO 22301 et ISO 31000 : continuité et risques}

ISO~22301 formule les exigences d'un système de management de la
continuité d'activité et est certifiable. ISO~31000 donne des lignes
directrices générales sur le management du risque, applicables à tout
type de risque, et n'est pas certifiable.

% ============================================================================
\section{Architecture, Projets et Maturité}
\label{sec:c04-projets}
% ============================================================================

\subsection{TOGAF : l'architecture d'entreprise}

TOGAF, publié par The Open Group, est le cadre de référence le plus
répandu en \termecle{architecture d'entreprise}\index{architecture
d'entreprise}~: la description cohérente des processus métier, des
données, des applications et des technologies d'une organisation, et de
leur évolution. Son cœur est une méthode de développement de
l'architecture (\textit{Architecture Development Method}, \sigle{ADM}).
TOGAF répond à la deuxième décision de Weill et Ross~: l'architecture
(Chapitre~\ref{chap:alignement}).

\subsection{PMBOK et PRINCE2 : la gestion de projet}

Le \textit{Guide to the Project Management Body of Knowledge}
(\sigle{PMBOK}), publié par le Project Management Institute, et la
méthode PRINCE2, d'origine britannique, sont les deux grandes références
de la gestion de projet. Le premier est un recueil de connaissances et de
principes~; la seconde, une méthode structurée en processus et en rôles.
Ils outillent l'objectif COBIT BAI11 (gérer les projets).

\subsection{CMMI : la maturité des processus}

Le CMMI (\textit{Capability Maturity Model Integration}), né à
l'université Carnegie Mellon pour évaluer les fournisseurs de logiciels du
département de la Défense américain, est aujourd'hui géré par l'ISACA. Il
décrit des niveaux de maturité, de~1 (processus improvisés) à~5
(processus en amélioration continue). COBIT~2019 s'en inspire directement
pour son modèle de capacité des processus (Chapitre~\ref{chap:cobit-capacite}).

% ============================================================================
\section{À Chaque Besoin son Référentiel}
\label{sec:c04-besoins}
% ============================================================================

Le choix d'un référentiel ne part pas du référentiel, mais du besoin.

\begin{tableaucours}{L{4.4cm}L{3.8cm}L{4.0cm}}{%
    \tableauHeader{Besoin de l'organisation} &
    \tableauHeader{Référentiel principal} &
    \tableauHeader{Compléments}
}
Organiser la gouvernance du conseil & ISO/IEC 38500, COBIT (EDM) & King~IV,
principes OCDE \\
Couvrir l'ensemble du management des SI & COBIT 2019 & · \\
Fiabiliser le contrôle interne & COSO (contrôle interne) & COBIT (MEA02) \\
Gérer les risques de l'entreprise & COSO ERM, ISO 31000 & COBIT (EDM03,
APO12) \\
Gérer les services au quotidien & ITIL & ISO/IEC 20000-1 \\
Faire reconnaître la qualité des services & ISO/IEC 20000-1 & ITIL \\
Sécuriser l'information & ISO/IEC 27001 et 27002 & COBIT (APO13, DSS05) \\
Assurer la continuité & ISO 22301 & ITIL, COBIT (DSS04) \\
Cohérence des systèmes & TOGAF & COBIT (APO03) \\
Conduire les projets & PMBOK, PRINCE2 & COBIT (BAI01, BAI11) \\
Améliorer la maturité des processus & CMMI & COBIT (capacité) \\
\end{tableaucours}
\vspace{-0.8em}
\begin{center}{\small\itshape Tableau~4.2. Correspondance entre besoins et
référentiels}\end{center}

\begin{boxobjet}{APO01 · Gérer le cadre de management des SI}
\textbf{Finalité~:} concevoir le système de management de l'information
et des technologies de l'organisation, et en garantir la cohérence avec
ses objectifs.

\textbf{Pratiques (extrait)~:} concevoir le système de management~;
communiquer les objectifs et l'orientation~; définir la structure
organisationnelle~; définir et entretenir les politiques et procédures.

\textbf{Sorties~:} politiques et procédures, définition des rôles,
\emph{choix et articulation des référentiels retenus}.

\textbf{RACI (extrait)~:} \emph{A} au directeur général, \emph{R} au DSI,
\emph{C} au responsable GRC et au RSSI.

\textbf{Indicateurs~:} part des politiques à jour et approuvées~;
couverture des processus par un référentiel identifié.
\end{boxobjet}
{\small\itshape Intitulé traduit par nos soins à partir de
COBIT~2019~\cite{isaca2018cobitobjectifs}.}

% ============================================================================
\section{La Carte d'Articulation}
\label{sec:c04-carte}
% ============================================================================

\subsection{Trois niveaux}

Les référentiels s'articulent naturellement en trois niveaux. Au
sommet, ceux qui s'adressent à l'organe de gouvernance. Au milieu, le
cadre qui organise l'ensemble du management des SI. À la base, les
référentiels spécialisés, chacun sur son domaine.

\begin{center}
\begin{tikzpicture}[
    niv/.style = {
        draw=CouleurPrimaire, fill=CouleurPrimaire, rounded corners=4pt,
        minimum width=12.6cm, minimum height=0.9cm, align=center,
        font=\footnotesize\bfseries\color{white}
    },
    mil/.style = {
        draw=CouleurBordDef, fill=CouleurDef, rounded corners=4pt,
        minimum width=12.6cm, minimum height=0.9cm, align=center,
        font=\footnotesize\bfseries\color{CouleurPrimaire}
    },
    spe/.style = {
        draw=CouleurBordDef, fill=white, rounded corners=3pt,
        text width=1.4cm, minimum height=1.3cm, align=center,
        anchor=north, inner sep=2pt,
        font=\tiny\bfseries\color{CouleurPrimaire}
    }
]
\node[niv] (g) at (0,2.6) {Gouvernance~: ISO/IEC 38500 · COBIT (EDM) ·
    COSO · King~IV};
\node[mil] (m) at (0,1.3) {Management des SI~: COBIT 2019 (APO · BAI ·
    DSS · MEA)};
\node[spe] at (-5.4,0.6) {Services\\ ITIL\\ ISO 20000};
\node[spe] at (-3.6,0.6) {Sécurité\\ ISO 27001\\ ISO 27002};
\node[spe] at (-1.8,0.6) {Continuité\\ ISO 22301};
\node[spe] at (0,0.6)    {Risques\\ ISO 31000\\ COSO ERM};
\node[spe] at (1.8,0.6)  {Architecture\\ TOGAF};
\node[spe] at (3.6,0.6)  {Projets\\ PMBOK\\ PRINCE2};
\node[spe] at (5.4,0.6)  {Maturité\\ CMMI};
\end{tikzpicture}
\end{center}
\vspace{-0.4em}
\begin{center}{\small\itshape Figure~4.1. La carte d'articulation des
référentiels}\end{center}

La règle d'articulation est simple~: \textbf{COBIT fixe les objectifs, les
référentiels spécialisés fournissent les pratiques}. COBIT~2019 le prévoit
lui-même~: pour chacun de ses objectifs, il indique les normes et
référentiels qui le détaillent.

\subsection{Construire la carte d'une organisation}

\begin{boxscenario}{Opération SANAGA. La carte de la banque}
La responsable GRC de la BANQUE SANAGA propose au conseil la carte
suivante, à partir des défaillances du prologue et des exigences de la
COBAC.
\begin{enumerate}
    \item \textbf{Gouvernance}~: principes de l'ISO/IEC~38500 dans la
          charte du comité stratégique des SI, objectifs EDM de COBIT.
    \item \textbf{Management}~: COBIT~2019 comme cadre unique, avec une
          évaluation de capacité sur les objectifs prioritaires
          (Partie~\ref{part:cobit}).
    \item \textbf{Services}~: les pratiques ITIL de gestion des
          changements, des mises en production et des incidents, en
          priorité.
    \item \textbf{Continuité}~: un plan conforme à ISO~22301, sans viser la
          certification dans un premier temps.
    \item \textbf{Sécurité}~: un SMSI conforme à ISO/IEC~27001, avec une
          certification visée à trois ans sur le périmètre des systèmes de
          paiement.
\end{enumerate}
\textbf{Ce qui est écarté, et pourquoi.} Le programme CMMI proposé par un
cabinet n'est pas retenu~: le modèle de capacité de COBIT suffit au
besoin de mesure, et la banque ne développe pas elle-même ses logiciels.
\end{boxscenario}

\begin{boxprocedure}
\textbf{La carte des référentiels, pièce de gouvernance.} Pour qu'un
auditeur puisse établir que le choix des référentiels est maîtrisé, il
doit trouver~:
\begin{itemize}
    \item un document qui liste les référentiels retenus, leur rôle et leur
          édition de référence~;
    \item la décision qui l'approuve, et sa date~;
    \item pour chaque norme certifiable visée, le périmètre et l'échéance~;
    \item la liste des référentiels écartés, avec la raison~;
    \item une date de revue, pour suivre les nouvelles éditions.
\end{itemize}
\end{boxprocedure}

\begin{boxmetier}{Responsable GRC}
\textbf{Choisir les référentiels sans les empiler}
\begin{enumerate}
    \item Partir des besoins et des obligations (lois, régulateur, clients),
          jamais de l'offre des cabinets.
    \item Retenir un cadre unique de gouvernance et de management, puis un
          référentiel spécialisé par domaine.
    \item Ne viser une certification que si un tiers l'exige ou si elle
          apporte un avantage identifié, et toujours sur un périmètre
          défini.
    \item Suivre les éditions~: une politique rédigée sur une édition
          périmée expose à des écarts lors de l'audit.
\end{enumerate}
\end{boxmetier}

\begin{boiteRemarque}{Les référentiels importés sont-ils adaptés à l'Afrique~?}
Les référentiels internationaux ont été conçus pour des organisations
dotées de moyens, de compétences et de marchés de prestataires que
beaucoup d'organisations camerounaises n'ont pas. Ce n'est pas une raison
de les écarter, mais d'en faire un usage \emph{proportionné}~: retenir les
objectifs essentiels, viser des niveaux de capacité réalistes, et
privilégier les pratiques qui traitent les risques locaux les plus
probables (dépendance aux prestataires, coupures d'énergie, rareté des
compétences). COBIT~2019 offre précisément cet outil d'adaptation avec
ses facteurs de conception (Chapitre~\ref{chap:cobit-conception}).
\end{boiteRemarque}

% ============================================================================
\separateur
% ============================================================================

\begin{boxresume}
\textbf{Ce qu'il faut retenir de ce chapitre~:}
\begin{itemize}
    \item \textbf{Trois natures}~: les cadres (COBIT, COSO) organisent, les
          normes certifiables (ISO/IEC~27001, 20000-1, ISO~22301) exigent,
          les recueils de bonnes pratiques (ITIL, PMBOK, ISO/IEC~27002)
          recommandent.
    \item \textbf{Certifier une organisation n'est pas certifier une
          personne}~: il n'existe pas de certification COBIT ou ITIL d'une
          entreprise.
    \item \textbf{Lire le périmètre}~: un certificat ne vaut que pour le
          périmètre qu'il désigne.
    \item \textbf{Trois niveaux}~: gouvernance (ISO/IEC~38500, COBIT EDM,
          COSO), management (COBIT~2019), domaines spécialisés (ITIL,
          ISO~27001, ISO~22301, TOGAF, PMBOK, CMMI\ldots).
    \item \textbf{COBIT fixe les objectifs, les autres fournissent les
          pratiques}~: ils ne se contredisent pas, ils se répartissent le
          travail.
    \item \textbf{Partir du besoin}~: une carte d'articulation se construit à
          partir des risques et des obligations, et se révise à chaque
          nouvelle édition.
    \item \textbf{La loi s'ajoute aux normes}~: au Cameroun, les audits de
          sécurité prévus par la loi n°~2010/012 coexistent avec les
          certifications volontaires.
\end{itemize}
\end{boxresume}

\bigskip

\begin{boxexo}[title={\ding{45}\quad Exercice~4.1. Quelle nature~? \niveaubase}]
Pour chacun des documents suivants, indiquez s'il s'agit d'un cadre de
référence, d'une norme certifiable ou d'un recueil de bonnes pratiques, et
si une organisation peut être certifiée sur sa base.
\begin{enumerate}[(a)]
    \item ISO/IEC~27001.
    \item ITIL.
    \item COSO ERM.
    \item ISO/IEC~20000-1.
    \item ISO~31000.
    \item COBIT~2019.
\end{enumerate}
\end{boxexo}

\begin{boxexo}[title={\ding{45}\quad Exercice~4.2. Lire un certificat \niveaumoyen}]
Un prestataire candidat à l'infogérance de KRIBILOG SA joint à son offre
un certificat ISO/IEC~27001 dont le périmètre est~: « Hébergement des
serveurs du centre de données de Casablanca ». Le service proposé à
KRIBILOG serait exploité depuis Douala.
\begin{enumerate}[(a)]
    \item Que prouve ce certificat, et que ne prouve-t-il pas pour le
          service proposé~?
    \item Quelles pièces complémentaires M.~MBARGA doit-il exiger~?
    \item Quelle clause proposez-vous d'insérer dans le contrat
          (Chapitre~\ref{chap:itil-niveaux})~?
\end{enumerate}
\end{boxexo}

\begin{boxexo}[title={\ding{45}\quad Exercice~4.3. La carte de NOUN PAY \niveauexpert}]
NOUN PAY, jeune entreprise de paiement mobile de Bafoussam, compte
45~salariés, héberge tout son SI chez un fournisseur d'informatique en
nuage étranger et utilise un modèle d'intelligence artificielle pour
noter les demandes de microcrédit. Ses investisseurs exigent « une
gouvernance des SI au niveau des standards internationaux ».
\begin{enumerate}[(a)]
    \item Identifiez ses besoins et ses obligations prioritaires.
    \item Construisez sa carte d'articulation sur le modèle de la
          Figure~4.1, en justifiant chaque référentiel retenu.
    \item Indiquez au moins deux référentiels que vous écartez, et
          pourquoi.
    \item Proposez un calendrier sur trois ans, compatible avec la taille
          de l'entreprise.
\end{enumerate}
\end{boxexo}

\bigskip

\begin{boxinfo}
\textbf{Transition.} La Partie~\ref{part:fondements} a posé les fondements
de la gouvernance des SI~: ses origines, ses défaillances, ses modèles et
ses référentiels. La Partie~\ref{part:cobit} entre dans le référentiel
central de ce cours. Le Chapitre~\ref{chap:cobit-archi} présente les
principes et l'architecture de COBIT~2019.
\end{boxinfo}
